\begin{pastexam}[label={px:gk3-098-ds-02}]{高考三級資訊處理 98 年 第二題}
（申論；25 分）

解釋下列名詞：（每小題 5 分共 25 分）
\begin{itemize}[leftmargin=2.2em]
\item[㈠] AVL 樹（AVL tree）
\item[㈡] 解釋圖形（graph）名詞：漢米爾頓迴路（Hamiltonian circuit）
\item[㈢] 解釋圖形（graph）名詞：廣度優先搜尋（breadth first search）。以程式實作此搜尋時，該使用那一種
資料結構？
\item[㈣] C++或 JAVA 語言中，protected 之意義
\item[㈤] Hanoi towers problem
\end{itemize}
\end{pastexam}

\begin{pastsoln}
五個名詞分屬不同章：㈠ 第~10~章，㈡ 第~28~章，㈢ 第~21~章，㈣ 是程式語言的物件導向觀念（資結書沒有
專章），㈤ 第~6~章（資結書第~6.4~節用遞迴解河內塔）。 兩小題是圖的名詞，本題放在第~21~章。 名詞解釋除了定義，最好再附一句
性質或例子，閱卷時比較好給分。

\textbf{㈠ AVL 樹}：一種自我平衡的二元搜尋樹（Adelson-Velsky 與 Landis，1962）。 每個節點的
\textbf{平衡因子}＝左子樹高度 $-$ 右子樹高度，AVL 樹要求每個節點的平衡因子都是 $-1$、$0$ 或 $+1$。
插入或刪除後若某個節點的平衡因子變成 $\pm 2$，用 LL、RR、LR、RL 四種旋轉恢復（資結書第~10.2~節）。
$n$ 個節點的 AVL 樹高度是 $O(\log n)$（約 $1.44 \lg n$ 以內），所以搜尋、插入、刪除最壞都是
$O(\log n)$，不會像一般二元搜尋樹那樣退化成 $O(n)$。

\textbf{㈡ 漢米爾頓迴路}：圖 $G = (V, E)$ 中\textbf{經過每個頂點恰好一次}、最後回到起點的迴路（環），
也就是長度 $|V|$ 的簡單環。 有漢米爾頓迴路的圖叫漢米爾頓圖。 和尤拉迴路（經過每條\textbf{邊}恰好一次）
不同：尤拉迴路看度數就能判斷（連通且每個頂點度數為偶數），判斷一張圖有沒有漢米爾頓迴路卻是
NP 完備問題（資結書第~28~章），目前沒有已知的多項式時間演算法。 例如完全圖 $K_n$（$n \ge 3$）一定有，
樹一定沒有。

\textbf{㈢ 廣度優先搜尋}：從起點 $s$ 出發，先拜訪 $s$ 的所有鄰居，再拜訪鄰居的鄰居，一層一層往外擴，
同一層（離 $s$ 邊數相同）的頂點都拜訪完才往下一層（資結書第~21.3~節）。 實作用\textbf{佇列}
（queue）：$s$ 標記後放進佇列；每次從前端取出一個頂點 $u$，把 $u$ 還沒標記的鄰居標記後放到尾端，直到
佇列空了為止。 佇列先進先出，保證離 $s$ 近的頂點先被處理。 用鄰接串列時間 $O(|V| + |E|)$，在無權圖上
順便求出 $s$ 到每個頂點的最少邊數。 （深度優先搜尋則用堆疊或遞迴。）

\textbf{㈣ protected}：類別成員的存取等級之一，介於 public 和 private 之間。 C++ 中，protected 成員
可以被\textbf{該類別自己的成員函式、friend，以及衍生類別（子類別）的成員函式}存取，類別外面的一般
程式碼不能存取。 Java 中，protected 成員可以被同一個 package 裡的類別，以及其他 package 中的子類別
存取。 用途是把實作細節對外隱藏（封裝），同時讓子類別在繼承時可以使用或改寫這些成員。

\textbf{㈤ 河內塔問題}：有三根柱子，第一根上有 $n$ 個由下往上越來越小的圓盤，要把全部圓盤搬到第三根，
規則是一次只能搬一個（某根柱子最上面的那個），且大盤不能放在小盤上。 遞迴解法（資結書虛擬碼~6.11）：先把上面 $n - 1$ 個
借第三根搬到第二根，把最大的一個搬到第三根，再把 $n - 1$ 個借第一根從第二根搬到第三根。
\begin{pseudo}[unbreakable]{河內塔}
\Procedure{Hanoi}{$n, \mathit{from}, \mathit{via}, \mathit{to}$}
    \If{$n \ge 1$}
        \State \Call{Hanoi}{$n - 1, \mathit{from}, \mathit{to}, \mathit{via}$}
        \State 把第 $n$ 個圓盤從 $\mathit{from}$ 搬到 $\mathit{to}$
        \State \Call{Hanoi}{$n - 1, \mathit{via}, \mathit{from}, \mathit{to}$}
    \EndIf
\EndProcedure
\end{pseudo}
搬動次數 $T(n) = 2T(n - 1) + 1$、$T(0) = 0$，展開得 $T(n) = 2^n - 1$（資結書第~2.5~節）。 這也是最少
次數：最大的盤要搬，必須其他 $n - 1$ 個都在第二根，之後它們還要再全部搬回來，所以至少
$2T(n - 1) + 1$ 次。 時間是 $\Theta(2^n)$，常拿來當指數時間遞迴的例子。
\end{pastsoln}
